\section{Picos de tráfico, Multi-Cloud y degradación elegante}

En las grandes plataformas comunitarias, los picos suelen originarse en noticias de última hora, eventos deportivos o políticos, o enlaces externos, y es difícil predecir del todo cuándo y dónde ocurrirán. El objetivo de la infraestructura no es que todas las funciones mantengan la misma calidad ante cualquier pico. Se trata de proteger las rutas centrales de lectura y escritura, la identidad y las capacidades de seguridad, y de degradar después, de manera progresiva, las funciones costosas según su valor para el negocio. La experiencia de Reddit se resume en cuatro pasos: absorb, divert, degrade y recover.

\subsection{Ruta de resiliencia en cuatro pasos}

El primer paso absorbe los picos ordinarios mediante cachés, colas y margen de capacidad elástica. El segundo evita las fallas locales al distribuir el tráfico entre varias regiones o varias nubes. El tercero, cuando la capacidad no alcanza, pasa al modo de solo lectura o desactiva las rutas de alto costo. El cuarto restablece la escritura, reproduce las colas y verifica el estado. \term{Graceful degradation}, o degradación elegante, significa que, ante falta de recursos o fallas parciales, el sistema conserva sus capacidades más importantes en lugar de caerse por completo.

\begin{figure}[H]
\centering
\includegraphics[width=0.96\textwidth]{images/09-traffic-resilience.png}
\caption{Ruta de resiliencia del tráfico en una plataforma comunitaria: absorción, desvío, degradación y recuperación.}
\figsource{Redibujada a partir de los subtítulos de la entrevista, 40:00--43:20, `images/09-traffic-resilience.png`.}
\end{figure}

\begin{knowledgebox}{Lectura de la figura: el orden de degradación debe codificarse de antemano}
Si se discute qué funciones importan cuando el incidente ya está en curso, la presión y la información parcial dominan las decisiones. El sistema debe definir de antemano la prioridad de la ruta central de lectura, la ruta de escritura de publicaciones y comentarios, el inicio de sesión, la moderación de seguridad, la búsqueda, las recomendaciones y las notificaciones. El modo de solo lectura mantiene visible el contenido de la comunidad, pero también puede impedir que los moderadores atiendan un daño que se está propagando. Por eso las herramientas de seguridad no deben desactivarse en bloque junto con las funciones de escritura ordinarias.
\end{knowledgebox}

\subsection{Multi-cloud no es replicar un conjunto de máquinas}

Usar varias nubes reduce el riesgo de que falle un solo proveedor y el riesgo en la negociación con él. A cambio, aumenta el costo de la replicación de datos, la identidad, la red, la observabilidad y los simulacros. Una arquitectura multinube solo aporta capacidad de recuperación si el equipo puede redirigir el tráfico en un tiempo aceptable, confirmar la consistencia de los datos y hacer que los servicios críticos funcionen de manera independiente en el segundo entorno. Desplegar unos cuantos recursos en otra nube sin haber hecho nunca un simulacro se parece más a diversificar las compras que a una conmutación por error.

\begin{importantbox}{Nota de clase: los objetivos de recuperación deben expresarse como acciones del usuario}
«La segunda nube está disponible» no es lo bastante concreto. Un mejor objetivo sería este: los usuarios pueden leer las publicaciones populares, los moderadores pueden ocultar contenido peligroso, las sesiones iniciadas no se invalidan de forma masiva y los comentarios nuevos pueden ponerse en cola y escribirse después de la recuperación. Expresar el objetivo como acciones del usuario pone al descubierto las dependencias y también permite saber si los simulacros cubren de verdad el valor central.
\end{importantbox}

\begin{table}[H]
\centering
\begin{tabular}{L{0.19\textwidth}L{0.31\textwidth}L{0.4\textwidth}}
\toprule
Capacidad & Impacto en el usuario si falla & Estrategia de degradación \\
\midrule
Lectura de contenido & La plataforma pierde casi toda su utilidad & Caché en varias capas, instantáneas estáticas, solo lectura entre regiones \\
Publicaciones y comentarios & Se interrumpe la participación en tiempo real & Colas, limitación de tasa, escritura diferida \\
Herramientas de seguridad para moderadores & El daño no puede contenerse a tiempo & Capacidad dedicada reservada e interfaz mínima de moderación \\
Ordenamiento de recomendaciones & Disminuye la personalización & Volver a lo más popular de la comunidad o al orden cronológico \\
Búsqueda/resúmenes con IA & El descubrimiento de información se vuelve más lento & Desactivar los resúmenes y conservar la búsqueda por palabras clave \\
Notificaciones y análisis & Se admite cierto retraso & Procesamiento por lotes y reproducción después de la recuperación \\
\bottomrule
\end{tabular}
\caption{La degradación se diseña según el valor para el usuario, no según qué equipo es dueño de cada microservicio.}
\end{table}

\begin{warningbox}{La reparación de datos después de la recuperación también es importante}
Que el tráfico se haya redirigido con éxito no significa que el incidente haya terminado. Las colas pueden consumirse más de una vez, los contadores pueden desviarse, el orden de los votos y los árboles de comentarios pueden quedar inconsistentes por un tiempo, y los eventos de seguridad pueden acumularse durante la ventana de degradación. El proceso de recuperación requiere reconciliation, eliminación de duplicados, auditoría y compensación a los usuarios. No basta con ver que vuelve a subir la tasa de éxito de HTTP.
\end{warningbox}

\subsection{Resumen de la sección}

La resiliencia de una plataforma depende de la capacidad reservada, de una redirección del tráfico que pueda ensayarse en simulacros, de una degradación elegante ordenada por el valor para el usuario y de la verificación de los datos después de la recuperación. Multi-cloud solo aporta resiliencia cuando la aplicación, los datos y la identidad pueden cambiar de entorno de verdad. A continuación se analiza cómo la IA cambia los puntos de entrada del tráfico y la forma de las interfaces, y qué necesidades de la comunidad no podrán sustituirse con modelos de resumen.
